% Options for packages loaded elsewhere
% Options for packages loaded elsewhere
\PassOptionsToPackage{unicode}{hyperref}
\PassOptionsToPackage{hyphens}{url}
\PassOptionsToPackage{dvipsnames,svgnames,x11names}{xcolor}
%
\documentclass[
  letterpaper,
  DIV=11,
  numbers=noendperiod]{scrartcl}
\usepackage{xcolor}
\usepackage{amsmath,amssymb}
\setcounter{secnumdepth}{-\maxdimen} % remove section numbering
\usepackage{iftex}
\ifPDFTeX
  \usepackage[T1]{fontenc}
  \usepackage[utf8]{inputenc}
  \usepackage{textcomp} % provide euro and other symbols
\else % if luatex or xetex
  \usepackage{unicode-math} % this also loads fontspec
  \defaultfontfeatures{Scale=MatchLowercase}
  \defaultfontfeatures[\rmfamily]{Ligatures=TeX,Scale=1}
\fi
\usepackage{lmodern}
\ifPDFTeX\else
  % xetex/luatex font selection
\fi
% Use upquote if available, for straight quotes in verbatim environments
\IfFileExists{upquote.sty}{\usepackage{upquote}}{}
\IfFileExists{microtype.sty}{% use microtype if available
  \usepackage[]{microtype}
  \UseMicrotypeSet[protrusion]{basicmath} % disable protrusion for tt fonts
}{}
\makeatletter
\@ifundefined{KOMAClassName}{% if non-KOMA class
  \IfFileExists{parskip.sty}{%
    \usepackage{parskip}
  }{% else
    \setlength{\parindent}{0pt}
    \setlength{\parskip}{6pt plus 2pt minus 1pt}}
}{% if KOMA class
  \KOMAoptions{parskip=half}}
\makeatother
% Make \paragraph and \subparagraph free-standing
\makeatletter
\ifx\paragraph\undefined\else
  \let\oldparagraph\paragraph
  \renewcommand{\paragraph}{
    \@ifstar
      \xxxParagraphStar
      \xxxParagraphNoStar
  }
  \newcommand{\xxxParagraphStar}[1]{\oldparagraph*{#1}\mbox{}}
  \newcommand{\xxxParagraphNoStar}[1]{\oldparagraph{#1}\mbox{}}
\fi
\ifx\subparagraph\undefined\else
  \let\oldsubparagraph\subparagraph
  \renewcommand{\subparagraph}{
    \@ifstar
      \xxxSubParagraphStar
      \xxxSubParagraphNoStar
  }
  \newcommand{\xxxSubParagraphStar}[1]{\oldsubparagraph*{#1}\mbox{}}
  \newcommand{\xxxSubParagraphNoStar}[1]{\oldsubparagraph{#1}\mbox{}}
\fi
\makeatother


\usepackage{longtable,booktabs,array}
\usepackage{calc} % for calculating minipage widths
% Correct order of tables after \paragraph or \subparagraph
\usepackage{etoolbox}
\makeatletter
\patchcmd\longtable{\par}{\if@noskipsec\mbox{}\fi\par}{}{}
\makeatother
% Allow footnotes in longtable head/foot
\IfFileExists{footnotehyper.sty}{\usepackage{footnotehyper}}{\usepackage{footnote}}
\makesavenoteenv{longtable}
\usepackage{graphicx}
\makeatletter
\newsavebox\pandoc@box
\newcommand*\pandocbounded[1]{% scales image to fit in text height/width
  \sbox\pandoc@box{#1}%
  \Gscale@div\@tempa{\textheight}{\dimexpr\ht\pandoc@box+\dp\pandoc@box\relax}%
  \Gscale@div\@tempb{\linewidth}{\wd\pandoc@box}%
  \ifdim\@tempb\p@<\@tempa\p@\let\@tempa\@tempb\fi% select the smaller of both
  \ifdim\@tempa\p@<\p@\scalebox{\@tempa}{\usebox\pandoc@box}%
  \else\usebox{\pandoc@box}%
  \fi%
}
% Set default figure placement to htbp
\def\fps@figure{htbp}
\makeatother
\usepackage{svg}





\setlength{\emergencystretch}{3em} % prevent overfull lines

\providecommand{\tightlist}{%
  \setlength{\itemsep}{0pt}\setlength{\parskip}{0pt}}



 


\KOMAoption{captions}{tableheading}
\makeatletter
\@ifpackageloaded{caption}{}{\usepackage{caption}}
\AtBeginDocument{%
\ifdefined\contentsname
  \renewcommand*\contentsname{Table of contents}
\else
  \newcommand\contentsname{Table of contents}
\fi
\ifdefined\listfigurename
  \renewcommand*\listfigurename{List of Figures}
\else
  \newcommand\listfigurename{List of Figures}
\fi
\ifdefined\listtablename
  \renewcommand*\listtablename{List of Tables}
\else
  \newcommand\listtablename{List of Tables}
\fi
\ifdefined\figurename
  \renewcommand*\figurename{Figure}
\else
  \newcommand\figurename{Figure}
\fi
\ifdefined\tablename
  \renewcommand*\tablename{Table}
\else
  \newcommand\tablename{Table}
\fi
}
\@ifpackageloaded{float}{}{\usepackage{float}}
\floatstyle{ruled}
\@ifundefined{c@chapter}{\newfloat{codelisting}{h}{lop}}{\newfloat{codelisting}{h}{lop}[chapter]}
\floatname{codelisting}{Listing}
\newcommand*\listoflistings{\listof{codelisting}{List of Listings}}
\makeatother
\makeatletter
\makeatother
\makeatletter
\@ifpackageloaded{caption}{}{\usepackage{caption}}
\@ifpackageloaded{subcaption}{}{\usepackage{subcaption}}
\makeatother
\usepackage{bookmark}
\IfFileExists{xurl.sty}{\usepackage{xurl}}{} % add URL line breaks if available
\urlstyle{same}
\hypersetup{
  pdftitle={No Country for Old Men},
  pdfauthor={Ganon Evans},
  colorlinks=true,
  linkcolor={blue},
  filecolor={Maroon},
  citecolor={Blue},
  urlcolor={Blue},
  pdfcreator={LaTeX via pandoc}}


\title{No Country for Old Men}
\author{Ganon Evans}
\date{}
\begin{document}
\maketitle


\subsection{Introduction}\label{introduction}

Patrick Womack was a 50-year old inmate at Texas's Coffield prison. On
an August in 2023, Womack was found dead in his bunk with an internal
body temperature of over 106 degrees Fahrenheit
(https://www.documentcloud.org/documents/24990424-10-patrick-womack\_redacted/).
The previous day, the heat index outside reached 113 degrees Fahrenheit.
In his unairconditioned cell, Womack asked a guard for the chance to
take a cold shower, but the guard declined because staff shortages meant
there weren't enough people there to watch him
(https://apnews.com/us-news/prisons-crime-texas-general-news-483428d72fdbcbe724700ddf675e15f2).
Though an autopsy report also linked Womack's death to his medical
history, heat may have played a significant role.

As summers get hotter in Texas, many of its prisons lack proper AC for
the over 100,000 prisoners there. Given this issue and the existent lack
of medical care in the prison system, I was interested in how one group
in particular may fare: the elderly. Older people are more susceptible
to conditions like heat stroke, and could have chronic conditions made
worse by extreme heat. In light of this, I wanted to explore how the
Texas prison system handles older prisoners: how many are there? What do
they have in common with other demographic groups? How do facilities
differ in their demographics?

As a consistent benchmark, I decide to label 50 and older as the
demographic I'd be following for this study.

\subsection{A Profile: Who of the elderly is in Texas
prisons?}\label{a-profile-who-of-the-elderly-is-in-texas-prisons}

As of the Texas Department of Criminal Justice's August 2025 profile, of
the 134,303 inmates across 105 facilities, 35,320 of them are age 50 or
older. 5.9\% of those 50 and older are women, compared to 8.3\% among
the whole population. The charges for inmates vary in expected sentence
length, from things like DWI to sexual assault or murder.

\begin{figure}[H]

{\centering \pandocbounded{\includesvg[keepaspectratio]{../images/test.svg}}

}

\caption{test123}

\end{figure}%

By and large, the older someone is in prison, their average sentence
length and time served increases linearly as well. In the data, there
was only 200 people above age 80; as age increases, the number of
inmates dwindles less and less, but significantly so after that point.

FIGURE 2

20\% of Texas's prisoners are in some stage of parole review. The
highest proportion of those in parole review relative to the rest of
their age group is for the 30-40 and 40-50 groups, which are also the
highest population deciles in the system. The proportion dwindles more
as the deciles get older, until only a small amount of those over 70 are
in the review process at all. This means that the oldest groups

PAROLE REVIEW FIGURE

\subsection{Race graphs}\label{race-graphs}

The three largest racial groups among inmates are Black, Hispanic, and
White. Looking at the relationship between age and tme served all three
groups exhibit the same positive relationship. Yet, after the mid-20s in
terms of age, Black inmates on average have served longer terms, then
Hispanic inmates, then White inmates. This implies that many long-term
Black inmates may have been imprisoned younger in life and have spent
more time dealing with the resources allocated to them through the
prisoner system. A follow-up to this report could connect medical
outcomes among the different races to see if age and length in prison
has any correlation.

The density of inmates over age 50 by race takes on a similar shape,
with White inmates having the highest density, especially in terms of
people who reach age 70+.

\subsection{Facility Section}\label{facility-section}

Scatterplot

Prisons with high elderly populations are concentrated in areas around
Southeast Texas

\subsection{Conclusion}\label{conclusion}

As Texas's inmates grow older, their likelihood of leaving prisons
decreases, whether its from decreased involvement in the parole process
or longer sentence times. This means that the health consequences of
aging like routine medical care or regulating cooling systems to not
excacerbate chronic conditions falls on the prisons. Yet, prisoners are
often at the bottom of society's list of groups to care for, as
evidenced by routine funding cuts and a lack of resources. Nobody lives
forever, but does a prisoner forfeit an additional 10 years of their
life ontop of their sentence when they're left without something as
basic as air conditioning? For a basic necessity like cooling buildings
to avoid sweltering heat, I think not.




\end{document}
